\chapter*{Preface}
\addcontentsline{toc}{chapter}{Preface}
\markboth{PREFACE}{PREFACE}

This monograph develops a geometric approach to the regularity problem for
the three-dimensional incompressible Navier--Stokes equations, one of the
seven Clay Millennium Problems.  The question, in its simplest form, asks
whether smooth initial data of finite energy necessarily give rise to a
globally smooth solution, or whether singularities can form in finite time.
Despite decades of progress in harmonic analysis, partial differential
equations, and geometric measure theory, the problem remains open.

\medskip

The central idea pursued here is to define a connection on the
infinite-dimensional bundle of divergence-free velocity fields, constructed
from the Biot--Savart kernel and the Leray projection at the natural
regularity of Leray--Hopf weak solutions.  Classical functional analytic
approaches to the regularity problem rely on energy methods that, at a
critical juncture, require estimates beyond what the energy class
$L^{\infty}_{t}(L^{2}_{\sigma}) \cap L^{2}_{t}(H^{1}_{0})$ provides.
The connection formulation offers a different path: by encoding the
geometric content of the flow into the curvature of a well-defined
connection on the divergence-free bundle, one remains entirely within the
Leray--Hopf class throughout.  The programme then asks whether
incompressibility, combined with the curvature constraints imposed by this
connection, obstructs finite-time blowup.

\medskip

A distinguishing feature of this work is its three-track architecture.
Every chapter is developed simultaneously along three parallel lines:

\begin{itemize}
  \item \textbf{Formal verification (Lean 4 and Mathlib).}
    All definitions, theorem statements, and (where completed) proofs are
    formalised in the Lean 4 proof assistant, building on the Mathlib
    library.  Lean serves as the final arbiter of correctness: a claim
    is accepted only when the type-checker confirms it.  The current
    formalisation contains a catalogue of \texttt{sorry} obligations
    corresponding to known theorems whose Lean proofs are deferred;
    these are tracked explicitly and do not represent gaps in the
    mathematical argument.

  \item \textbf{Theory (LaTeX and SymPy).}
    The mathematical narrative is written as a self-contained book,
    with complete proofs for all results.  Symbolic computation in SymPy
    independently verifies key identities, exponents, and estimates,
    providing a computational cross-check on the analytic arguments.

  \item \textbf{Numerics.}
    Numerical simulations (planned for later chapters) will test the
    geometric predictions of the theory against direct computation of
    the Navier--Stokes equations, using spectral and finite-element
    solvers.  The numerical track is not a substitute for proof, but
    a source of intuition and a check against false conjectures.
\end{itemize}

At each chapter boundary, a formal sync point reconciles the three tracks:
Lean definitions must match LaTeX definitions, SymPy verifications must
confirm the identities stated in the text, and (where applicable) numerical
experiments must be consistent with the theoretical predictions.  Lean is
the authority on ties.

\medskip

\textbf{Status of the programme.}
This document records an ongoing research programme.  As of the current
writing, Chapter~1 (functional analytic foundations) is complete across all
three tracks: the Lean formalisation compiles, the LaTeX proofs are
self-contained, and the SymPy verifications pass.  The remaining chapters
are in various stages of development.  \emph{We do not claim to have
resolved the Millennium Problem.}  The approach may ultimately succeed or
fail; this text is an honest record of the attempt, written so that the
mathematical community can evaluate, critique, and build upon the ideas
presented here.

\bigskip

\hfill Alejandro Jos\'e Soto Franco

\hfill \today
